%%%%%%%%%%%%%%%%%%%%%%%%%%%%%%%%%%%%%%%%%
% Article Notes
% Structure Specification File
% Version 2.3 (16/01/25)
%
% File made by frcrav
%
%%%%%%%%%%%%%%%%%%%%%%%%%%%%%%%%%%%%%%%%%

%----------------------------------------------------------------------------------------
%	REQUIRED PACKAGES
%----------------------------------------------------------------------------------------

%\usepackage[includeheadfoot,
%columnsep=2cm,
%left=1in,
%right=1in,
%top=.5in,
%bottom=.5in]{geometry} % Margins

% % two col. geometry
% \usepackage[
% 	includeheadfoot,
% 	columnsep=3cm,
% 	left=1in,
% 	right=1in,
% 	top=.5in,
% 	bottom=.5in
% ]{geometry} % Margins

% single page geometry
\usepackage[
	includeheadfoot,
	left=2.5cm,
	right=2.5cm,
	top=2cm,
	bottom=2cm
]{geometry} % Margins

\usepackage[utf8]{inputenc} % Required for inputting international characters

\usepackage[T1]{fontenc} % For international characters

\usepackage[english]{babel} % Use your language to modify automatic texts. English by default

% Fonts: XCharter, cmbright (overall nice), sourcesanspro (nice contrast with math font), tgpagella (math font is not incredible)
\usepackage[default]{XCharter} % XCharter as the main font
%\usepackage[default]{sourcesanspro} % \usepackage[default]
\newcommand{\xcharter}{\fontfamily{XCharter-TLF}\selectfont} % Used in \lecture section;
\usepackage{sourcesanspro} % \usepackage[default]


%----- Secondary packages
%\usepackage{showframe} % Show page frame for debugging

\usepackage{natbib} % Use natbib to manage the reference
\bibliographystyle{apalike} % Citation style

\usepackage{titlesec}


\usepackage{graphicx} % Required for inserting images
\usepackage{wrapfig} % Wrap figures and tables around the text

\usepackage{amsmath, amsthm, amsfonts, amssymb, mathtools, eurosym}

\usepackage{tikz, pgfplots} % plots and figures
\pgfplotsset{compat=1.18}
\usetikzlibrary{shapes.geometric, arrows, positioning, backgrounds}

\usepackage{circuitikz} % electrical drawings

\usepackage{listings, lstautogobble}
\usepackage[dvipsnames]{xcolor}
\usepackage{algorithm}
\usepackage{algpseudocode}

\usepackage{hyperref, enumerate, enumitem, color}

\usepackage[acronym]{glossaries}

\usepackage{lipsum}

\usepackage{tabularx} % More table options
\usepackage{booktabs}

\usepackage{fancyhdr}

\usepackage{marginnote}

\usepackage{etoolbox}
\usepackage{changepage}

\usepackage{tcolorbox}
\tcbuselibrary{theorems, skins, most}

\usepackage{xparse}

\usepackage{caption, subcaption}

\usepackage{siunitx}

\usepackage{multicol}




%----------------------------------------------------------------------------------------
%	CUSTOM COMMANDS
%----------------------------------------------------------------------------------------

\newcommand{\coursetitle}[1]{\renewcommand{\coursetitle}{#1}} % Define a command for storing the course title
\newcommand{\doctitle}{\coursetitle} % Define a command to store the article information as it will appear in the title and header

\newcommand{\datenotesstarted}[1]{\renewcommand{\datenotesstarted}{#1}} % Define a command to store the date when notes were first made
\newcommand{\docdate}[1]{\renewcommand{\docdate}{#1}} % Define a command to store the date line in the title
\newcommand{\lesyear}[1]{\renewcommand{\lesyear}{#1}}

\newcommand{\docauthor}[1]{\renewcommand{\docauthor}{#1}} % Define a command for storing the article notes author

% Define a command for the structure of the document title
\newcommand{\printtitle}{
	\thispagestyle{empty}
	\begin{center}
		\textbf{\Huge{\doctitle}} % Style 1
		%       \textsc{\huge{\doctitle}} % Style 2
		%       \xcharter\textbf{\huge\textsc{\doctitle}} % Style 3, XCharter family

		\large{\docdate}

		\large{\docauthor}

	\end{center}

	\vfill
	\begin{center}
		%	\fontseries{sb}
		%	\selectfont
		Lecture notes \\
		Year {\lesyear}
		\smallskip
		\begin{figure}[h]
			\centering
			\includegraphics[width=.35\linewidth]{/Users/frcrav/Desktop/LaTeX/Logo_PoliTo.png}
		\end{figure}

		\smallskip
		Politecnico di Torino, Italy
	\end{center}
	\newpage
}

% ----- Intro

% \makeintro{year1/year2}{intro}
\newcommand{\makeintro}[2] {
	\bigskip
	\section*{Introduction}
	% Use this section to explain the course, exam modes and informations
	\emph{#2}

	\vfill

	\begin{center}
		\fontseries{sb}
		\selectfont
		Lecture notes \\
		Year #1
		\smallskip
		\begin{figure}[h]
			\centering
			\includegraphics[width=.35\linewidth]{/Users/frcrav/Desktop/LaTeX/Logo_PoliTo.png}
		\end{figure}

		\smallskip
		Politecnico di Torino, Italy
	\end{center}
}

\newcommand{\toc}{
	\newpage\pagenumbering{roman}\pagestyle{plain}
	\tableofcontents
	\newpage\pagenumbering{arabic}\pagestyle{fancy}
}

\let\oldquote\quote
\let\endoldquote\endquote
\renewenvironment{quote}[2][]
{\if\relax\detokenize{#1}\relax
		\def\quoteauthor{#2}%
	\else
		\def\quoteauthor{#2~---~#1}%
	\fi
	\oldquote}
{\par\nobreak\smallskip\hfill(\quoteauthor)%
	\endoldquote\addvspace{\bigskipamount}}

% ----- Math

%\newtheorem{definition}{Definizione}[section]
%\newtheorem*{definition*}{Definizione}
%\newtheorem{theorem}{Teorema}[section]
%\newtheorem*{theorem*}{Teorema}
%\newtheorem{lemma}{Lemma}[section]
%\newtheorem{example}{Esempio}[section]
%\newtheorem*{example*}{Esempio}

% Theorems and math-env
\theoremstyle{plain}
\newtheorem{theorem}{Theorem}[section]
\newtheorem*{theorem*}{Theorem}
\newtheorem{lemma}{Lemma}[section]

\theoremstyle{definition}
\newtheorem{definition}{Definition}[section]
\newtheorem*{definition*}{Definition}
\newtheorem{example}{Example}[section]
\newtheorem*{example*}{Example}
\newtheorem{exs}{Exercise}
\newtheorem*{exs*}{Exercise}

\tcbset{
	examplestyle/.style={
			sharp corners,
			colback = white,
			before skip = 0.2cm,
			after skip = 0.5cm,
			boxrule = .5pt,
			colframe = black,
			% Opzionale: aggiunge un po' di spazio interno
			left=2mm, right=2mm, top=1mm, bottom=1mm
		}
}
\tcolorboxenvironment{example}{
	examplestyle  %
}
\newcommand{\solution}{
	\tcblower           % Questo comando di tcolorbox divide il box in due
	\textbf{Solution.}~ % Scrive "Solution." all'inizio della parte bianca
}



% \tcbset{
% 	sharp corners,
% 	colback = white,
% 	before skip = 0.2cm,    % add extra space before the box
% 	after skip = 0.5cm      % add extra space after the box
% }
% \newtcolorbox{examplebox}{
% 	boxrule = .5pt,
% 	colframe = black % frame color
% }

\theoremstyle{remark}
\newtheorem*{remark}{Remark} % Osservazioni
\newtheorem{note}{Observation}	 % Note
\newtheorem*{note*}{Observation}	 % Note


\NewDocumentEnvironment{sol}{O{}}
{\par\vspace{1em}%
	\noindent\textbf{Solution}%
	\if\relax\detokenize{#1}\relax
		%
	\else
		~(\textit{qst. #1})
	\fi
	.\par}
{\par\vspace{1em}}

%\newlength{\mymargin} \setlength{\mymargin}{2.7mm}
%
%\AtBeginEnvironment{theorem}{\begin{adjustwidth}{\mymargin}{}}
%	\AtEndEnvironment{theorem}{ \end{adjustwidth}}
%
%\AtBeginEnvironment{theorem*}{\begin{adjustwidth}{\mymargin}{}}
%	\AtEndEnvironment{theorem*}{\end{adjustwidth}}
%
%\AtBeginEnvironment{definition}{\begin{adjustwidth}{\mymargin}{}}
%	\AtEndEnvironment{definition}{\end{adjustwidth}}
%\AtBeginEnvironment{definition*}{\begin{adjustwidth}{\mymargin}{}}
%	\AtEndEnvironment{definition*}{\end{adjustwidth}}
%
%\AtBeginEnvironment{example}{\begin{adjustwidth}{\mymargin}{}}
%	\AtEndEnvironment{example}{\end{adjustwidth}}
%\AtBeginEnvironment{example*}{\begin{adjustwidth}{\mymargin}{}}
%	\AtEndEnvironment{example*}{\end{adjustwidth}}
%
%\AtBeginEnvironment{lemma}{\begin{adjustwidth}{\mymargin}{}}
%	\AtEndEnvironment{lemma}{\end{adjustwidth}}
%
%\AtBeginEnvironment{exs}{\begin{adjustwidth}{\mymargin}{}}
%	\AtEndEnvironment{exs}{\end{adjustwidth}}
%\AtBeginEnvironment{exs*}{\begin{adjustwidth}{\mymargin}{}}
%	\AtEndEnvironment{exs*}{\end{adjustwidth}}

% Math mode symbols for common sets and spaces. Example usage: $\R$
\newcommand{\R}{\mathbb{R}}	% Real numbers
\newcommand{\C}{\mathbb{C}}	% Complex numbers
\newcommand{\Q}{\mathbb{Q}}	% Rational numbers
\newcommand{\Z}{\mathbb{Z}}	% Integers
\newcommand{\N}{\mathbb{N}}	% Natural numbers
\newcommand{\F}{\mathcal{F}}	% Calligraphic F for a sigma algebra
\newcommand{\El}{\mathcal{L}}	% Calligraphic L, e.g. for L^p spaces

% Math mode symbols for probability
\newcommand{\pr}{\mathbb{P}}	% Probability measure
\newcommand{\E}{\mathbb{E}}	% Expectation, e.g. $\E(X)$
\newcommand{\var}{\text{Var}}	% Variance, e.g. $\var(X)$
\newcommand{\cov}{\text{Cov}}	% Covariance, e.g. $\cov(X,Y)$
\newcommand{\corr}{\text{Corr}}	% Correlation, e.g. $\corr(X,Y)$
\newcommand{\B}{\mathcal{B}}	% Borel sigma-algebra

% Other miscellaneous symbols
\newcommand{\tth}{\text{th}}	% Non-italicized 'th', e.g. $n^\tth$
\newcommand{\Oh}{\mathcal{O}}	% Big-O notation, e.g. $\O(n)$

\DeclareMathOperator{\length}{length}
\DeclareMathOperator{\Aut}{Aut}
\DeclareMathOperator{\diam}{diam}
\DeclareMathOperator*{\res}{res}

% ----- Code, algo

% ----- Colors and highlighter
\definecolor{myGreen}{RGB}{48, 209, 88}
\definecolor{myGreenDark}{RGB}{0, 137, 50}

\newcommand{\emphcustom}[2]{{\color{#1}\textit{#2}}} % Custom color on emph text
\newcommand{\emphtb}[1]{{\textbf{\textit{#1}}}} % Boldface emph
\newcommand{\emphred}[1]{{\color{red}\textit{#1}}} % Red emph
\newcommand{\emphblue}[1]{{\color{blue}\textit{#1}}} % Blue emph
\newcommand{\emphgreen}[1]{{\color{myGreenDark}\textit{#1}}} % Green emph

\newcommand{\emphr}[1]{{\sffamily\color{red}\MakeUppercase{\textbf{#1}}}} % Red emph


\newcommand{\hl}[1]{\colorbox{yellow}{#1}} % Yellow highlither
\newcommand{\hlb}[1]{\colorbox{yellow}{\textbf{#1}}} % Yellow highlither with bold text

% ----- Other cmds



%----------------------------------------------------------------------------------------
%	STRUCTURE MODIFICATIONS
%----------------------------------------------------------------------------------------

% === Custom Lecture Command ===
\newcounter{lecture} % Create lecture counter
\renewcommand{\thelecture}{\ifnum\value{lecture}<10 0\fi\arabic{lecture}}

\newcommand{\lecture}[2]{%
	\clearpage
	\refstepcounter{lecture} % increase lecture number automatically
	%\refstepcounter{lecture} % Lable referencing

	\section*{\xcharter Les \thelecture: #1 \\ \small #2}% prints the title unnumbered
	\addcontentsline{toc}{section}{Les \thelecture: #1 (#2)}% adds it to the table of contents

	\setcounter{section}{0}% reset section numbering
	\setcounter{subsection}{0}% reset subsection numbering
	\setcounter{subsubsection}{0}% reset subsubsection numbering
}

\newcounter{assignment} % Create lecture counter
\renewcommand{\theassignment}{\ifnum\value{assignment}<10 0\fi\arabic{assignment}}
\newcounter{question} % Create lecture counter

\newcommand{\assignment}[2]{%
	\newpage
	\refstepcounter{assignment} % increase lecture number automatically

	\section*{\xcharter Ass \theassignment: #1 \\ \small #2}% prints the title unnumbered
	\addcontentsline{toc}{section}{Ass \theassignment: #1 (#2)}% adds it to the table of contents

	\setcounter{section}{0}% reset section numbering
	\setcounter{subsection}{0}% reset subsection numbering
	\setcounter{subsubsection}{0}% reset subsubsection numbering
	\setcounter{question}{0}
}

% --- Definizione ambiente "Question" ---
\newtcbtheorem[use counter=question]{qst}{Question}%
{colback=white,
	colframe=black!40,
	fonttitle=\bfseries,
	coltitle=black,
	boxrule=0.4pt,
	sharp corners,
	enhanced,
	attach boxed title to top left={yshift=-2mm, xshift=2mm},
	boxed title style={colback=white, boxrule=0.4pt, colframe=black!40}}%
{qst}

\setlength{\parskip}{3pt} % Slightly increase spacing between paragraphs
\setlength{\parindent}{0pt} % Eliminate paragraph indentation.

\makeatletter
\renewcommand\paragraph{\@startsection{paragraph}{4}{\z@}%
	{3.25ex \@plus1ex \@minus.2ex}%
	{-1em}%
	{\normalfont\normalsize\fontseries{b}\selectfont\underline}}
\makeatother

\makeatletter
\renewcommand\subparagraph{\@startsection{subparagraph}{5}{\parindent}
	{3.25ex \@plus1ex \@minus .2ex}
	{-1em}
	{\normalfont\normalsize\underline}}
\makeatother
% Uncomment to center section titles
%\usepackage{sectsty}
%\sectionfont{\centering}

% Uncomment for Roman numerals for section numbers
%\renewcommand\thesection{\Roman{section}}

\pagestyle{fancy} % Custom fancy header
\fancyhf{}
\setlength{\headheight}{14pt} % For single page
\fancyhead[RE,LO]{\MakeUppercase \doctitle}
\fancyhead[RO,LE]{\thepage}
\fancyfoot[C]{Les \thelecture}
